% !TEX root = ../main_hessenberg.tex
\section{AI collaboration record}\label{sec:appendix-ai}

Following the transparency conventions proposed in~\cite{schmitt2025}, we
document the role of AI systems in producing the formalization described in
Section~\ref{sec:methodology}. Throughout this appendix, human-authored
paragraphs are marked with a blue bar and AI-generated text with an orange bar.

\subsection{Agent architecture}

\begin{humanparagraph}
The formalization was produced using the following system:
\begin{itemize}
\item \textbf{Claude Code} (Opus~4.6): the primary Lean code generator, with
  access to file read/write, \texttt{lake build}, and LSP diagnostics.
\item \textbf{lean-prover sub-agent} (Aristotle~\cite{achim2025aristotleimolevelautomatedtheorem}): a
  specialized agent with tools for checking tactic states, trying multiple
  proof strategies in parallel, and parsing compiler errors.
\item \textbf{Human}: strategy design, goal decomposition, and debugging.
\end{itemize}
\end{humanparagraph}

\subsection{Proof difficulty classification}

\begin{aiparagraph}
The table below classifies each major theorem by difficulty and the degree of
human intervention required.

\begin{center}
\renewcommand{\arraystretch}{1.2}
\footnotesize
\begin{tabular}{@{}p{5.4cm}ccl@{}}
\toprule
\textbf{Theorem} & \textbf{Difficulty} & \textbf{Autonomy} & \textbf{Human role} \\
\midrule
\texttt{SignEquiv.refl/symm} & Low & Full & --- \\
\texttt{loop\_at\_first/last/mid} & Low & Full & --- \\
\texttt{hcycle\_subset} & Low & Full & --- \\
\texttt{backward\_arc\_is\_spine} & Low & Full & --- \\
\texttt{cut\_lemma} & Low & Full & --- \\
\texttt{loopless\_iff} & Medium & Full & --- \\
\texttt{outDeg\_first}, \texttt{inDeg\_last} & Medium & High & Set decomposition hint \\
\texttt{givensRotation\_orthogonal} & Medium & High & Pythagorean identity hint \\
\texttt{givensProduct\_isUpperHessenberg} & Medium & Partial & \texttt{partialGivensProduct\_apply\_of\_row\_ge} helper \\
\texttt{inDeg\_mem\_S} & Medium & Partial & Set decomposition \\
\texttt{completeness} & Medium & High & $\pi/4$ vs.\ $\pi/2$ construction \\
\texttt{count\_indep\_subsets} & Medium & High & Partition strategy \\
\texttt{fixes\_last\_implies\_id} & High & Partial & Downward induction \\
\texttt{digraph\_model} & High & Low & 5-lemma decomposition \\
\texttt{singleton\_iso} & High & Partial & Modular arithmetic \\
\texttt{singleton\_iso\_unique} & High & Low & In-degree characterization \\
\texttt{rigid} & Highest & Low & Cut/pigeonhole strategy \\
\bottomrule
\end{tabular}
\end{center}
\end{aiparagraph}

\subsection{Case study: the rigidity proof}

\begin{humanparagraph}
The proof of \texttt{rigid} (Theorem~\ref{thm:rigid}) was the single most
challenging component, requiring multiple rounds of strategy revision.

\medskip
\noindent\textbf{Attempt~1: Degree argument.}
The initial idea was to show that vertex~$1$ is the unique vertex with
out-degree $|S| + 1$, forcing any isomorphism to fix it. This fails because
for singletons $|S| = 1$, another vertex can share the same out-degree.

\medskip
\noindent\textbf{Attempt~2: Chain along Hamilton cycle.}
If $\sigma$ fixes vertex~$n$, backward spine arcs force
$\sigma(n{-}1) = n{-}1$, $\sigma(n{-}2) = n{-}2$, etc. This reasoning is
correct and became \texttt{fixes\_last\_implies\_id}, but it requires first
proving that $\sigma$ \emph{does} fix vertex~$n$.

\medskip
\noindent\textbf{Attempt~3: Cut/pigeonhole (successful).}
The human proposed:
\begin{enumerate}
\item Suppose $\sigma(n) = w \neq n$.
\item In-degree analysis forces $w + 1 \in S'$.
\item The cut lemma constrains arcs crossing the cut at $w$.
\item Case split on $\sigma(1)$ relative to $w$:
  \begin{itemize}
  \item $\sigma(1) < w$: pigeonhole gives a contradiction
    (\texttt{perm\_image\_below\_cut}).
  \item $\sigma(1) > w$: out-degree counting gives a contradiction
    (\texttt{perm\_image\_above\_cut}).
  \end{itemize}
\item Hence $\sigma$ fixes $n$; apply \texttt{fixes\_last\_implies\_id}.
\end{enumerate}
Each sub-lemma was stated by the human with explicit hypotheses; the AI agent
proved each independently. The final assembly was straightforward.
\end{humanparagraph}

\subsection{Case study: the bridge theorem}

\begin{humanparagraph}
The bridge theorem (\texttt{digraph\_model}) was decomposed into five lemmas:
\begin{enumerate}
\item \texttt{givensProduct\_row\_vanishing}: $\cos(\theta_{i-1}) = 0$ implies row~$i$ is
  zero in the upper triangle.
\item \texttt{givensProduct\_col\_vanishing}: $\cos(\theta_j) = 0$ implies column~$j$ is
  zero in the upper triangle.
\item \texttt{partialGivensProduct\_propagation}: $\cos(\theta_{i-1}) \neq 0$ implies entry
  $(i, m{+}1)$ is nonzero for all $m \ge i$.
\item \texttt{partialGivensProduct\_apply\_of\_col\_ge}: columns beyond the first $m$ Givens factors are
  untouched.
\item \texttt{givensProduct\_apply\_subdiag\_ne\_zero}: subdiagonal entries equal
  $\sin(\theta_j) \neq 0$.
\end{enumerate}
The agent initially tried to prove \texttt{digraph\_model} monolithically and
failed. After the human provided the five-lemma decomposition, each piece was
proved independently. The introduction of
\texttt{givensRotation.apply\_of\_ne\_col} (a reusable API for Givens entry
evaluation) resolved repeated failures from ad-hoc \texttt{givensRotation}
unfolding.
\end{humanparagraph}

\subsection{Failure modes and lessons}

\begin{humanparagraph}
\begin{enumerate}
\item \textbf{Lost work via \texttt{git checkout}} (twice). Uncommitted proofs
  were discarded. \emph{Lesson}: commit after every successful proof.

\item \textbf{LSP timeouts}. Complex lemmas caused the Lean server to time
  out, which the agent misinterpreted as proof failure.
  \emph{Lesson}: set \texttt{maxHeartbeats} proactively.

\item \textbf{Unproductive unfolding}. The agent repeatedly expanded
  \texttt{givensRotation} instead of using named lemmas.
  \emph{Lesson}: encapsulate entry-level facts into reusable API lemmas.

\item \textbf{Tactic explosion in modular arithmetic}. The proofs of
  \texttt{singleton\_iso} and \texttt{singleton\_iso\_unique} generated
  hundreds of subgoals from modular arithmetic case splits.
  \emph{Lesson}: use \texttt{omega} and \texttt{grind} aggressively.

\item \textbf{Strategy persistence}. When a proof approach failed, the agent
  sometimes persisted with minor variations rather than reconsidering.
  \emph{Lesson}: set retry limits and escalate to human review.
\end{enumerate}
\end{humanparagraph}

\subsection{Representative strategy prompts}

\begin{humanparagraph}
\noindent\textbf{Prompt~1} (Rigidity via cut/pigeonhole):
\begin{quote}\itshape
Suppose $\sigma(n{-}1) \neq n{-}1$, so $\sigma$ maps vertex~$n$ to some
$w < n$. Use in-degree analysis to show $w{+}1 \in S'$, then the cut lemma to
constrain all values. Split on whether $\sigma(0) < w$ or $\sigma(0) \ge w$.
Below: pigeonhole. Above: out-degree contradiction.
\end{quote}

\noindent\textbf{Prompt~2} (Bridge decomposition):
\begin{quote}\itshape
Don't prove \texttt{digraph\_model} directly. First prove: (1) row vanishing,
(2) col vanishing, (3) propagation, (4) identity cols, (5) subdiag nonzero.
Each by induction on $m$ using the \texttt{List.ofFn} splitting.
\end{quote}

\noindent\textbf{Prompt~3} (Completeness construction):
\begin{quote}\itshape
Set $\theta_k = \pi/4$ if $k{+}1 \in S$, else $\pi/2$.
Then $\sin \neq 0$ always and $\cos \neq 0 \iff k{+}1 \in S$.
Use \texttt{positivity} for $\sqrt{2}/2 \neq 0$.
\end{quote}
\end{humanparagraph}
